\section{Little Big Last Man}

\begin{quotex}
The Constitution of 1795, like its predecessors, was made for man. But there is no such thing as man in the world. In my lifetime I have seen Frenchmen, Italians, Russians, etc.; thanks to Montesquieu, I even know that one can be Persian. But as for man, I declare that I have never in my life met him; if he exists, he is unknown to me.
\flright{\textsc{Joseph de Maistre}}
\end{quotex}


In the 1970s film, ``Little Big Man'', the Cheyenne tribes referred to themselves as ``human beings'', as all other peoples were aliens. But isn't that the traditional view? For the Greeks, non-Greeks were ``barbarians''. The Romans had a similar view.

It is only the modern Western man that has decided that we are all ``human beings''. Yet the modern's notion of ``human being'' is just as parochial as the Cheyenne's. That is because the ``human being'' looks remarkably like the ideal of the Enlightenment European: universal, cosmopolitan, rational, egalitarian, materialist, scientific, ``inclusive''. Blinded by his own ideology and convinced he now knows man qua man, he is unable to notice any differences between different peoples.

Modern European man believes that, having thrown off all conventions of the past -- for he views all religions and social customs as mere conventions -- he has discovered man. What he has discovered is Nietzsche's ``last man''. This is the result of successive declines beginning with the Reformation and Enlightenment, continuing with liberal capitalism and socialism, and finally concluding with globalisation. The latter is the most insidious since it is a softer and more deceptive version of the previous economic systems.

The great theoretician of globalisation was \textbf{Alexandre Kojève}. Initially a Marxist, he soon realised that global capitalism would be more effective in bringing about the egalitarian state; he was an early participant in the World Trade Organisation, the cornerstone of globalisation.

In his book \emph{An Introduction to the Reading of Hegel}, Kojève foresaw the end of man as such through his re-animalisation. With all needs satisfied and conventions overthrown, there would be no need for human discourse:

\begin{quotex}
The definitive annihilation of Man properly so-called also means the definitive disappearance of human Discourse (Logos) in the strict sense. Animals of the species Homo sapiens would react by conditioned reflexes to vocal signals of sign language, and thus their so-called discourses would be like what is supposed to be the ``language'' of bees''. What would disappear, then, is not only Philosophy or the search for discursive Wisdom, but also that Wisdom itself. For in these post-historical animals, there would no longer be any discursive understanding of the World and of self.
\end{quotex}


Kojève came to this conclusion after his visits to the United States and realised that the Soviets and Chinese were simply poorer versions of Americans. But he believed the ``American way of life'' was the prefiguration for the rest of the world. Undifferentiated, classless, a ``melting pot'' of races and cultures, religion merely a private affair, feminist, sexually ``liberated'' --- this way of life was what Kojève not only saw, but actively promoted. Rarely does a philosopher get the opportunity to promote his views as a matter of public policy.

The last man is materially satiated and liberated from the conventions and superstitions of the past. So he has nothing left to talk about and his vocal sounds are as insignificant as the chirping of crickets. This discovery of ``man'' by the ruling elites of Europe and the United States explains the otherwise inexplicable policies they implement.

This all stems from --- as Kojève indicates --- the denial of the ``Logos''; that is, it is just one more revolutionary movement, in a string of such movements, that seeks to overturn every established order in favour of an undifferentiated, promiscuous, immanentist society.

Contrast this with the traditional society based on the ``Logos'': transcendent, spiritual, differentiated, hierarchical. It sees its social/political structure and religion as essential elements, not as mere conventions. There is, and has been, a deliberate, even if unconscious, effort to overturn all such societies.


\flrightit{Posted on 2007-12-16 by Cologero}

\begin{center}* * *\end{center}

\begin{footnotesize}
\begin{sffamily}
\texttt{Logres on 2012-01-28 at 10:52 said:}

George Parkin Grant (a Canadian) did some wonderful critiques of Kojeve. This is a fine article, very succinct.


\hfill

\end{sffamily}
\end{footnotesize}

